\chapter{Análisis Diferencial con el Sistema de Referencia}

El análisis forense digital alcanza su mayor nivel de certidumbre probatoria cuando se fundamenta en el **análisis diferencial**: la comparación sistemática y reproducible entre la imagen del sistema comprometido y una instalación de referencia limpia, idéntica en versión de kernel, distribución y paquetes de software base.

En este capítulo se expone la metodología empleada en el laboratorio sobre la estación \textit{Dell Precision 350}, detallando el análisis binario estático y la extracción de cadenas de texto (\textit{strings}) que permitieron diseccionar los binarios del sistema troyanizados por el atacante (\texttt{/bin/ps}, \texttt{/bin/netstat}, \texttt{/bin/ls} y \texttt{/sbin/init}).

\section{Metodología de Cotejo Diferencial en Laboratorio}

Para discriminar con exactitud las modificaciones maliciosas frente a archivos de configuración legítimos, el equipo pericial montó en el laboratorio de referencia dos discos duros bajo una instalación idéntica de Debian GNU/Linux:
\begin{itemize}
    \item \textbf{Volumen de Control (Limpio):} Contenía los binarios originales distribuidos por el empaquetado oficial del sistema operativo, verificados mediante sus sumas de verificación criptográficas originales.
    \item \textbf{Volumen Forense (Copia bit a bit):} Réplica montada en modo de solo lectura del almacenamiento de \texttt{victima1}.
\end{itemize}

Mediante el comando de verificación de paquetes y la ejecución automatizada de \texttt{chkrootkit}, se auditaron 1,999 binarios ejecutables presentes en las ramas \texttt{/bin}, \texttt{/sbin}, \texttt{/usr/bin} y \texttt{/usr/sbin}. El análisis identificó que los binarios esenciales de administración de procesos y red habían sido reemplazados por variantes pertenecientes a la familia clásica del \textit{Linux Rootkit} (LRK).

\section{Análisis de /bin/ps: Filtrado de Procesos y Puerta Trasera}

El comando \texttt{/bin/ps} es la herramienta primordial utilizada por los administradores de sistemas para interrogar la tabla de procesos del kernel a través del sistema de archivos virtual \texttt{/proc}. La versión localizada en el servidor comprometido presentaba una discrepancia sustancial en su tamaño en disco y en su firma hash MD5 respecto al binario de control.

Al someter el archivo a un análisis de cadenas de texto legibles (\texttt{strings -a /bin/ps}), se recuperaron los artefactos documentados en el Listado~\ref{lst:ps_strings}.

\begin{lstlisting}[style=terminal,caption={Extracto de cadenas forenses extraídas del binario troyanizado /bin/ps.},label={lst:ps_strings}]
/lib/ld-linux.so.2
libc.so.6
strcpy
fgets
system
malloc
setuid
setgid
$Id: ps.c,v 1.11 2001/09/03 Exp $
[1;33m%s
/bin/su -
/dev/.libpt
\end{lstlisting}

\subsection{Evidencia del Cuerpo del Delito en /bin/ps}

El examen técnico de las cadenas extraídas reveló de forma contundente la mecánica de la troyanización:
\begin{enumerate}
    \item \textbf{Filtrado activo mediante \texttt{/dev/.libpt}:} La presencia de la ruta codificada en duro \texttt{/dev/.libpt} demuestra que el binario fue recompilado para abrir y leer dicho archivo al iniciar su ejecución. Cada vez que el administrador ejecutaba \texttt{ps}, el programa comparaba el nombre de cada proceso en memoria contra la lista negra de \texttt{.libpt} (\texttt{dnsquery}, \texttt{sysinit}, \texttt{popclean}), omitiendo deliberadamente su impresión en pantalla si coincidían.
    \item \textbf{Puerta trasera de ejecución (\textit{backdoor}):} La inclusión de la cadena \texttt{/bin/su -} y la llamada a \texttt{system()} y \texttt{setuid()} confirman que el binario integraba una compuerta encubierta. Al invocar \texttt{ps} con un modificador no documentado o bajo una variable de entorno secreta, el programa ejecutaba un intérprete de comandos con privilegios de superusuario sin requerir contraseña.
    \item \textbf{Firma del kit de desarrollo:} La etiqueta RCS \texttt{\$Id: ps.c,v 1.11 2001/09/03 Exp \$} vinculó formalmente el binario con el código fuente del kit \textit{Linux Rootkit versión 5} (LRK5), ampliamente distribuido en foros clandestinos en la época.
\end{enumerate}

\section{Análisis de /bin/netstat: Evasión de Conexiones de Red}

De manera análoga a \texttt{ps}, el binario \texttt{/bin/netstat} fue modificado para manipular la visualización de los sockets de red abiertos en el sistema operativo.

El análisis de cadenas del binario reveló referencias cruzadas hacia el archivo de configuración \texttt{/dev/.libn}:

\begin{lstlisting}[style=terminal,caption={Cadenas identificadas en el binario modificado /bin/netstat.},label={lst:netstat_strings}]
/lib/ld-linux.so.2
libc.so.6
fopen
fgets
fclose
strstr
/dev/.libn
6667
17749
56321
\end{lstlisting}

Al ejecutarse, este \texttt{netstat} adulterado interceptaba la lectura de las tablas \texttt{/proc/net/tcp} y \texttt{/proc/net/udp}. Si una conexión activa involucraba los puertos definidos en \texttt{.libn} (como el puerto IRC 6667 o los puertos de retransmisión 17749 y 56321), el programa descartaba silenciosamente el socket antes de formatear la salida en terminal. De esta manera, el administrador que supervisaba las conexiones de red permanecía en la creencia errónea de que el servidor no mantenía ninguna comunicación con el exterior.

\section{Compendio Comparativo de Binarios Comprometidos}

El Cuadro~\ref{tab:comparativa_binarios} consolida las discrepancias criptográficas y funcionales detectadas entre los binarios originales de control y los ejecutables troyanizados hallados en el sistema comprometido.

\begin{table}[htbp]
\centering
\small
\begin{tabular}{lp{4cm}p{4cm}p{3cm}}
\toprule
\textbf{Binario Auditado} & \textbf{Hash MD5 (Control)} & \textbf{Hash MD5 (Comprometido)} & \textbf{Mecanismo de Evasión} \\
\midrule
\texttt{/bin/ps} & \texttt{e4d8cb98...} (64~KB) & \texttt{9a12bc44...} (88~KB) & Filtra procesos según \texttt{/dev/.libpt}; backdoor \texttt{su -}. \\
\texttt{/bin/netstat} & \texttt{7b21fa09...} (92~KB) & \texttt{3c88de11...} (112~KB) & Oculta puertos 6667 y sockets de bots IRC. \\
\texttt{/bin/ls} & \texttt{1f44aa72...} (48~KB) & \texttt{6e90ff32...} (68~KB) & Oculta entradas con prefijos especiales en \texttt{/dev}. \\
\texttt{/sbin/killall5} & \texttt{55cc9011...} (18~KB) & \texttt{88dd4100...} (24~KB) & Evita matar los procesos de los bots al apagar. \\
\texttt{/usr/bin/top} & \texttt{fa091122...} (52~KB) & \texttt{12aa5677...} (74~KB) & Oculta el consumo de CPU de los procesos del atacante. \\
\bottomrule
\end{tabular}
\caption{Comparación criptográfica y funcional entre binarios limpios y troyanizados.}
\label{tab:comparativa_binarios}
\end{table}

\section{Conclusiones del Análisis Diferencial}

El análisis diferencial en el laboratorio forense aportó certezas definitivas para la investigación:
\begin{itemize}
    \item La alteración no fue producto de una infección viral ciega o aleatoria, sino de una **sustitución deliberada y coordinada** de las herramientas de administración del sistema operativo.
    \item Los binarios troyanizados formaban un ecosistema cohesivo: \texttt{ps} y \texttt{top} ocultaban la ejecución de los bots, \texttt{netstat} enmascaraba el tráfico hacia los servidores de IRC, y \texttt{ls} velaba los archivos de configuración en el disco duro.
    \item La fecha y marcas de inodo de estos reemplazos suministraron los anclajes temporales indispensables para la siguiente fase pericial: la construcción de la línea de tiempo global del sistema.
\end{itemize}
